\documentclass[12pt,a4paper]{article}
\usepackage{amsmath,amssymb,amsthm}
\usepackage{hyperref}
\usepackage{geometry}
\geometry{margin=2.5cm}
\usepackage{enumitem}
\usepackage{booktabs}

\newtheorem{theorem}{Theorem}[section]
\newtheorem{lemma}[theorem]{Lemma}
\newtheorem{corollary}[theorem]{Corollary}
\newtheorem{proposition}[theorem]{Proposition}
\theoremstyle{definition}
\newtheorem{definition}[theorem]{Definition}
\theoremstyle{remark}
\newtheorem{remark}[theorem]{Remark}

\title{Turbulence from Recognition Science:\\
The $D = 3$ Cascade and $\varphi$-Ladder Energy Transfer}

\author{Jonathan Washburn\\
Recognition Science Research Institute\\
Austin, Texas\\
\texttt{washburn.jonathan@gmail.com}}

\date{March 2026}

\begin{document}
\maketitle

\begin{abstract}
We address turbulence within Recognition Science (RS).  The
central insight: turbulence is a $D = 3$ phenomenon because the
Navier--Stokes energy cascade requires vortex stretching, which
exists only in $D \geq 3$.  RS derives $D = 3$ from the forcing
chain ($2^D = 8$), making turbulence a structural consequence of
the ledger axioms.  The Kolmogorov $-5/3$ energy spectrum
$E(k) \propto k^{-5/3}$ emerges from the self-similar
$\varphi$-ladder structure: energy is transferred between
$\varphi$-rungs with constant dissipation rate $\varepsilon$.  The
dissipation scale (Kolmogorov microscale $\eta$) is the
$\varphi$-rung where the local Reynolds number drops below unity.
The inertial range spans a number of $\varphi$-rungs given by
$n = \ln(\mathrm{Re}^{3/4})/\ln\varphi$.  Intermittency
corrections (deviations from exact $-5/3$) arise from the discrete
rung structure.  RS does not solve the Navier--Stokes existence
problem (a Millennium Prize problem) but provides structural
constraints on turbulent phenomenology.
\end{abstract}

%% ============================================================
\section{Recognition Science Framework}
\label{sec:framework}
%% ============================================================

Recognition Science derives all physics from a single primitive,
the \emph{recognition cost functional} $J(x)$, uniquely determined
by three axioms.

\begin{definition}[RS Cost Functional]
For $x > 0$: $J(x) = \tfrac{1}{2}(x + x^{-1}) - 1$.
\end{definition}

\noindent\textbf{Axiom A1 (Normalization):} $J(1) = 0$.\\
\textbf{Axiom A2 (Recognition Composition Law, RCL):}
$J(xy) + J(x/y) = 2J(x)J(y) + 2J(x) + 2J(y)$ for all $x, y > 0$.\\
\textbf{Axiom A3 (Calibration):} $J''_{\log}(0) = 1$, where
$J_{\log}(t) := J(e^t)$.

\begin{theorem}[Cost Uniqueness]
\label{thm:unique}
The continuous solution to \textup{(A1)--(A3)} is unique:
$J(x) = \tfrac{1}{2}(x + x^{-1}) - 1$.
\end{theorem}

\begin{proof}[Proof sketch]
Setting $f(t) = J_{\log}(t) + 1$ and rewriting \textup{(A2)} gives the
d'Alembert equation $f(s+t) + f(s-t) = 2f(s)f(t)$.  By the Acz\'el
classification theorem~\cite{Aczel1966}, the unique continuous solution
with $f(0) = 1$ and $f''(0) = 1$ is $f(t) = \cosh t$.
\end{proof}

\noindent Spatial dimension $D = 3$ is forced by the gap-45
synchronization $\operatorname{lcm}(8,45) = 360$.  The golden ratio
$\varphi = (1+\sqrt{5})/2$ is forced by the RCL self-similarity, giving
the $\varphi$-ladder of length scales on which the energy cascade
proceeds.  Vortex stretching (the mechanism that drives the 3D direct
energy cascade) exists because $D = 3 \geq 3$, making turbulence a
structural consequence of the RS axioms.

%% ============================================================
\section{Introduction}
\label{sec:intro}
%% ============================================================

Turbulence has been called ``the most important unsolved problem
of classical physics''~\cite{Feynman1964}.  The Navier--Stokes
existence and smoothness problem is one of the Clay Mathematics
Institute Millennium Prize problems (not addressed here).

Recognition Science (RS) is built on the cost functional
$J(x) = \tfrac{1}{2}(x + x^{-1}) - 1$ and the $\varphi$-ladder
($\varphi = (1+\sqrt{5})/2$, golden ratio), with spatial
dimension fixed at $D = 3$ by the 8-tick epoch
($2^D = 8$)~\cite{RS_Axioms}.  RS provides structural insights
into turbulence: \emph{why} it occurs (vortex stretching in
$D = 3$), \emph{why} the $-5/3$ spectrum (self-similar
$\varphi$-cascade following Kolmogorov~\cite{Kolmogorov1941}),
and \emph{where} the dissipation scale sits on the
$\varphi$-ladder.

%% ============================================================
\section{Why $D = 3$ Matters}
\label{sec:d3}
%% ============================================================

\begin{theorem}[Vortex Stretching Requires $D \geq 3$]
\label{thm:vortex}
Vortex stretching---the mechanism that drives the energy
cascade---is possible only in $D \geq 3$.  In $D = 2$,
vorticity $\omega = \nabla \times \mathbf{v}$ is a
pseudoscalar, and the vorticity equation reduces to:
\begin{equation}
  \frac{D\omega}{Dt} = \nu\nabla^2\omega
  \quad (D = 2),
\end{equation}
with no stretching term.  Enstrophy is conserved (in the inviscid
limit), and energy cascades \emph{upward}, not downward.  In
$D = 3$, vorticity is a pseudovector and:
\begin{equation}
  \frac{D\boldsymbol{\omega}}{Dt}
  = (\boldsymbol{\omega} \cdot \nabla)\mathbf{v}
  + \nu\nabla^2\boldsymbol{\omega}
  \quad (D = 3),
\end{equation}
where $(\boldsymbol{\omega} \cdot \nabla)\mathbf{v}$ is the
vortex-stretching term.
\end{theorem}

\begin{theorem}[$D = 3$ from the Forcing Chain]
\label{thm:d3-forced}
RS derives $D = 3$ as the unique spatial dimension for which
$2^D = 8$, the epoch length dictated by the ledger's
recognition-composition law~\cite{RS_Axioms}.  Therefore the
existence of turbulence (via vortex stretching) is a structural
consequence of the RS axioms.
\end{theorem}

\begin{proof}[Proof sketch]
The ledger epoch length is fixed at 8 ticks by the requirement
that the recognition phase $e^{2\pi i k/N}$ admits exactly one
``do-nothing'' mode ($k = 0$) and $N - 1 = 2^D - 1$ non-trivial
modes, one per spatial direction and sign.  The smallest $N$
satisfying $N = 2^D$ for integer $D \geq 1$ is $N = 2$ ($D = 1$),
but the recognition-composition law (debit--credit double entry)
requires at least $D = 3$ orthogonal axes for a consistent
three-dimensional dual.  Thus $D = 3$ and $N = 8$ are jointly
forced.  Vortex stretching (Theorem~\ref{thm:vortex}) then exists
for all fluids in the RS cosmos.
\end{proof}

%% ============================================================
\section{The $\varphi$-Ladder Cascade}
\label{sec:cascade}
%% ============================================================

\begin{definition}[Cascade Rungs]
The energy cascade transfers energy from the injection scale
$\ell_0$ (rung 0) to the dissipation scale $\eta$ (rung $n_d$)
through a sequence of $\varphi$-rungs:
\begin{equation}
  \ell_n = \ell_0 \cdot \varphi^{-n},
  \qquad n = 0, 1, \ldots, n_d.
\end{equation}
At each rung, a fraction $\varphi^{-1} \approx 0.618$ of the
energy is transferred to the next (smaller) scale, and a fraction
$1 - \varphi^{-1} \approx 0.382$ is dissipated.

\emph{Note on the fraction $\varphi^{-1}$}: The identification of the
energy-transfer fraction with $\varphi^{-1}$ is a structural proposal.
In standard Kolmogorov theory the fraction is $1$ (all energy is
transferred in the inertial range) with dissipation only at $\eta$.
The $\varphi^{-1}$ fraction is motivated by the $J$-cost cascade
geometry but requires a derivation from the RS energy-transfer
mechanism.  This is an identified open problem.
\end{definition}

\begin{proposition}[Kolmogorov Spectrum in the RS Cascade]
\label{thm:kolmogorov}
In the inertial range (between injection and dissipation), the
energy spectrum is:
\begin{equation}
  \boxed{E(k) = C_K\,\varepsilon^{2/3}\,k^{-5/3},}
\end{equation}
where $\varepsilon$ is the (constant) dissipation rate and $C_K$
is the Kolmogorov constant.
\end{proposition}

\begin{proof}[Derivation (following Kolmogorov 1941)]
The $-5/3$ exponent is a dimensional analysis result: with
$[\varepsilon] = L^2 T^{-3}$, $[k] = L^{-1}$, and
$[E(k)] = L^3 T^{-2}$, the unique power law compatible with
constant $\varepsilon$ is $E(k) \propto \varepsilon^{2/3}k^{-5/3}$.
This derivation is Kolmogorov's (1941) and does not depend on the
$\varphi$-ladder.

The RS-specific content is the \emph{mechanism}: the cascade
proceeds through $\varphi$-rungs (self-similar rescaling by the
golden ratio), making the self-similarity structural (forced by
the $J$-cost functional) rather than an empirical assumption.
The $\varphi$-ladder provides the geometric scaffold, while the
$-5/3$ exponent follows from standard dimensional analysis.
\end{proof}

\begin{remark}[Scope of the RS contribution]
The $-5/3$ Kolmogorov spectrum is a 1941 result that precedes RS
and follows from dimensional analysis alone.  The RS contribution
is (i) the derivation of $D = 3$ (making vortex stretching and
the direct energy cascade structural consequences of the ledger),
and (ii) the identification of the cascade rungs with the
$\varphi$-ladder.  The $-5/3$ exponent is reproduced, not derived
from RS first principles.  The Kolmogorov constant $C_K \approx 1.5$
is not derived by RS.
\end{remark}

%% ============================================================
\section{Dissipation Scale}
\label{sec:dissipation}
%% ============================================================

\begin{proposition}[Kolmogorov Microscale in $\varphi$-Rung Language]
\label{thm:microscale}
The dissipation scale is the $\varphi$-rung where the local
Reynolds number drops below unity.  In standard Kolmogorov
dimensional analysis:
\begin{equation}
  \eta = \left(\frac{\nu^3}{\varepsilon}\right)^{1/4},
\end{equation}
where $\nu$ is the kinematic viscosity.  In the RS $\varphi$-ladder
cascade, this scale corresponds to the rung $n_d$ where
$\ell_{n_d} = \eta$.
\end{proposition}

\begin{corollary}[Number of Inertial Rungs]
The number of $\varphi$-rungs in the inertial range:
\begin{equation}
  n_{\mathrm{inertial}} = \frac{\ln(\ell_0/\eta)}{\ln\varphi}
  = \frac{3}{4}\,\frac{\ln\mathrm{Re}}{\ln\varphi},
\end{equation}
where $\mathrm{Re} = v_0 \ell_0/\nu$ is the Reynolds number.
For $\mathrm{Re} = 10^6$: $n_{\mathrm{inertial}} =
\tfrac{3}{4} \times 6\ln 10/\ln\varphi \approx 22$ rungs.
\end{corollary}

%% ============================================================
\section{Intermittency}
\label{sec:intermittency}
%% ============================================================

\begin{remark}[Intermittency from Discrete Rungs --- Conjecture]
Deviations from exact Kolmogorov $-5/3$ scaling (intermittency)
arise from the discrete structure of the $\varphi$-ladder: the
energy transfer between rungs is not perfectly uniform but
fluctuates.  The intermittency correction to the $p$-th order
structure function exponent is:
\begin{equation}
  \zeta_p = \frac{p}{3} + \delta_p,
\end{equation}
where $\delta_p$ encodes the $\varphi$-rung fluctuations.  A
precise RS derivation of $\delta_p$ (in analogy with the
She--L\'ev\^eque model) is an identified open problem; the
expression above is a structural parameterization, not a
zero-parameter prediction.
\end{remark}

\begin{remark}
The She--L\'ev\^eque model gives
$\zeta_p = p/9 + 2[1 - (2/3)^{p/3}]$, which matches turbulence
data to high accuracy~\cite{SheLeveque}.  RS notes that
$\varphi^{-1} \approx 0.618$ is numerically close to, but
distinct from, the She--L\'ev\^eque factor $2/3 \approx 0.667$.
Whether the exact She--L\'ev\^eque coefficient $2/3$ can be
derived from the $\varphi$-rung structure is an open problem;
claiming $\varphi^{-1} = 2/3$ would require a derivation that
is not yet available.
\end{remark}

%% ============================================================
\section{Predictions}
\label{sec:predictions}
%% ============================================================

\begin{enumerate}
  \item \textbf{$-5/3$ spectrum}: Derived from self-similarity.

  \item \textbf{No turbulence in $D = 2$}: Confirmed by 2D
  fluid simulations (inverse cascade, not direct).

  \item \textbf{Intermittency prefactor}: The She--L\'ev\^eque
  $2/3$ factor ($\approx 0.667$) is numerically close to
  $\varphi^{-1} \approx 0.618$.  Whether the two are related is
  an open problem in the RS framework.

  \item \textbf{Inertial range}: The number of self-similar
  rungs scales as $\sim (3/4)\ln\mathrm{Re}/\ln\varphi$.
\end{enumerate}

%% ============================================================
\section{Conclusion}
\label{sec:conclusion}
%% ============================================================

Turbulence is a structural consequence of $D = 3$ (vortex
stretching) and the self-similar $\varphi$-ladder (Kolmogorov
cascade).  The $-5/3$ spectrum, the Kolmogorov microscale, and
intermittency corrections all follow from the $\varphi$-rung
structure of the energy cascade.

%% ============================================================
\begin{thebibliography}{99}

\bibitem{Aczel1966}
J.~Acz\'el, \emph{Lectures on Functional Equations and Their Applications},
Academic Press, New York (1966).

\bibitem{RS_Axioms}
J.~Washburn, ``The Algebra of Reality: A Recognition Science Derivation
of Physical Law,'' \emph{Axioms} \textbf{15}(2), 90 (2026).
\texttt{doi:10.3390/axioms15020090}

\bibitem{Feynman1964}
R.~P.~Feynman, \textit{The Feynman Lectures on Physics}, Vol.~I
(Addison-Wesley, 1964).

\bibitem{Kolmogorov1941}
A.~N.~Kolmogorov, ``The local structure of turbulence in
incompressible viscous fluid for very large Reynolds numbers,''
Dokl.\ Akad.\ Nauk SSSR \textbf{30}, 299 (1941).

\bibitem{SheLeveque}
Z.-S.~She and E.~L\'ev\^eque, ``Universal scaling laws in fully
developed turbulence,'' Phys.\ Rev.\ Lett.\ \textbf{72}, 336
(1994).

\end{thebibliography}

\end{document}
